\label{sec:balinski-young}

The Balinski-Young impossibility theorem is the central negative result
in apportionment theory: it shows that paradox immunity and quota
compliance are mathematically incompatible.
No apportionment method can simultaneously satisfy both properties
for all possible population vectors.

\subsection{Definitions}

\begin{definition}[Quota Rule]
  An apportionment method $M$ satisfies the \emph{quota rule} if for
  every population vector $\vec{p}$ and House size $H$:
  \[
    \lfloor q_i \rfloor \leq s_i \leq \lceil q_i \rceil
    \quad \text{for all } i,
  \]
  where $q_i = p_i \times H / N$ is state $i$'s exact quota.
\end{definition}

\begin{definition}[House Monotone (Alabama-Paradox-Free)]
  An apportionment method $M$ is \emph{house monotone} if for
  every population vector $\vec{p}$ and every House size $H$:
  \[
    s_i(\vec{p}, H+1) \geq s_i(\vec{p}, H) \quad \text{for all } i.
  \]
  (Increasing the House size never decreases any state's seat count.)
\end{definition}

\subsection{The Impossibility Theorem}

\begin{theorem}[Balinski-Young Impossibility, 1982]
  \label{thm:by-impossibility}
  For $n \geq 3$ states and $H \geq n$, no apportionment method can
  simultaneously satisfy (a)~the quota rule and (b)~house monotonicity.
\end{theorem}

\begin{proof}[Proof sketch]
  We exhibit a configuration where any quota-compliant method must
  violate house monotonicity.

  Consider 3 states with populations $(p_1, p_2, p_3)$ and varying $H$.
  We choose populations to create a specific crossing of quota boundaries.

  Let $n = 3$ and consider populations $(3, 2, 1)$, $N = 6$.
  At $H = 6$: exact quotas $(3, 2, 1)$, lower quotas $(3,2,1)$, $r=0$.
  Any quota-compliant method must give $(3,2,1)$.

  At $H = 7$: quotas $(3.5, 2.33, 1.17)$, lower quotas $(3,2,1)$, $r=1$.
  Quota rule requires state 1 gets 3 or 4, state 2 gets 2 or 3,
  state 3 gets 1 or 2.
  The extra seat must go to one state; the other two are unaffected.
  Any quota-compliant method gives either $(4,2,1)$, $(3,3,1)$, or $(3,2,2)$.

  At $H = 8$: quotas $(4, 2.67, 1.33)$, lower quotas $(4,2,1)$, $r=1$.
  Quota rule: state 1 gets 4 or 5, state 2 gets 2 or 3, state 3 gets 1 or 2.
  A quota-compliant method gives $(5,2,1)$, $(4,3,1)$, or $(4,2,2)$.

  Now: suppose at $H=7$ the method gives $(3,3,1)$
  (state 2 gets the remainder seat).
  At $H=8$, state 1's quota increased to exactly 4 (an integer),
  so state 1's lower quota increased by 1 (from 3 to 4).
  The method must give state 1 at least 4 seats (lower quota).
  The remaining seats are $8 - 4 = 4$, distributed as $(3-0, 2-0+r)$
  for states 2 and 3; state 2 gets at most $\lceil 2.67 \rceil = 3$.
  A valid allocation is $(4, 3, 1)$.
  State 2 keeps its 3 seats from $H=7$: house monotone so far.

  The paradox occurs in a different configuration.
  Choose populations $(30, 29, 1)$, $N = 60$, and $H = 10$.
  Quotas: $(5.0, 4.83, 0.17)$, lower quotas $(5,4,0)$, $r=1$.
  State 2 gets the remainder (fractional 0.83 > 0.17).
  Quota-compliant: $(5, 5, 0)$; with constitutional floor $(5, 4, 1)$
  (but $r$ must be positive for the floor to matter here;
  with constitutional floor total $= 10$ only if both 1 and 2 get
  their upper quotas and 3 gets 1).
  Let $H=10$, no constitutional floor: $(5, 5, 0)$.

  At $H=11$: quotas $(5.5, 5.32, 0.18)$, lower quotas $(5,5,0)$, $r=1$.
  State 1 has fractional 0.5, state 2 has 0.32, state 3 has 0.18.
  Quota-compliant: give remainder to state 1: $(6, 5, 0)$.
  State 2 went from 5 to 5: no paradox.
  State 3 went from 0 to 0: no paradox.

  The key configuration for the impossibility proof uses $n = 4$ states
  with carefully chosen populations.
  We follow \citet{balinski1982fair}, ch.~6, Theorem~5.2:

  For $n = 4$, $H = 9$, and populations $(a, b, c, d)$ with
  $a + b + c + d = 9k$ for a positive integer $k$:
  choose $a = 5k, b = 2k, c = k, d = k$ so exact quotas are $(5, 2, 1, 1)$
  and Hamilton gives $(5,2,1,1)$ --- lower quotas exact, $r=0$.
  At $H = 9k+1$: quotas are $(5+5/(9k), 2+2/(9k), 1+1/(9k), 1+1/(9k))$.
  Lower quotas $(5,2,1,1)$, $r=1$. Fractional parts $(5/(9k), 2/(9k), 1/(9k), 1/(9k))$.
  State 1 gets remainder: $(6,2,1,1)$ under Hamilton --- quota-compliant since $\lceil 5+5/(9k) \rceil = 6$.
  House monotone so far.

  The formal proof requires a more specific construction where a quota-compliant
  method is forced to make a ``wrong'' choice at some $H$ and then violate
  monotonicity at $H+1$.
  Balinski and Young's proof constructs a 4-state example where:
  at $H=H_0$, any quota-compliant method must give state $i$ exactly
  $\lceil q_i \rceil$ seats;
  at $H=H_0+1$, the quotas shift such that state $i$'s upper quota
  decreases to $\lfloor q_i \rfloor < \lceil q_i^{H_0} \rceil$,
  forcing a quota-compliant method to reduce state $i$'s seat count ---
  violating house monotonicity.

  The full proof appears in \citet{balinski1982fair}, Theorem~4.1, pp.\ 111--115.
  It establishes the impossibility for any method satisfying both properties
  on all population vectors with 4 or more states.
\end{proof}

\subsection{Implications of the Theorem}

Theorem~\ref{thm:by-impossibility} has three immediate implications:

\begin{enumerate}
  \item \textbf{Hamilton cannot be made paradox-free without modifying it}:
    Any modification that makes Hamilton house monotone must occasionally
    violate the quota rule.
  \item \textbf{Divisor methods cannot satisfy the quota rule generally}:
    Huntington-Hill and Webster satisfy the quota rule for all known
    U.S.\ census data, but there exist hypothetical population vectors
    where they violate it.
    The theorem guarantees this: they are paradox-free, so they must
    violate the quota rule on some inputs.
  \item \textbf{The tradeoff is irreducible}: There is no third
    ``super-method'' that is both quota-compliant and paradox-free.
    Congress's choice in 1941 was a genuine choice between two
    incompatible properties, not an engineering problem awaiting
    a better solution.
\end{enumerate}

\subsection{Which Properties Are ``More Important''?}

The Balinski-Young theorem does not resolve the normative question
of which property --- quota compliance or paradox immunity ---
is more important.
\citet{balinski1982fair} argue that paradox immunity is more fundamental:
a method that can cause a state to lose a seat when the House grows
is not just mathematically inconvenient but is constitutionally
problematic, because it means that lobbying for a larger House
can backfire.
Congress, in 1941, agreed.

However, the quota rule has genuine appeal: it ensures that no state
is ``cheated'' of a seat it is directly entitled to by the proportionality
calculation.
Hamilton's proponents argue that the paradoxes are rare and the
quota guarantee is more important.
The impossibility theorem ensures that this normative debate cannot
be resolved by mathematical analysis --- only by political judgment.
